\chapter[Programmazione Lineare: geometria e ottimo]{Programmazione Lineare: forme, geometria e proprietà dell'ottimo}\label{cap:pl}

Parte del programma che inizia qui: risoluzione grafica, forme alternative, geometria della PL, proprietà della soluzione ottima; seguiranno il metodo del simplesso e la dualità. In tutto il capitolo le variabili sono \textbf{continue}.

\section{Il coltivatore, riscalato}

Esprimendo la resa in migliaia di euro e lo stallatico in decine di tonnellate, il modello di \S\ref{sec:colt} diventa
\[
\max\ 3x_L+5x_P\quad \st\ x_L+x_P\le12,\ \ x_L\le10,\ \ x_P\le6,\ \ x_L+2x_P\le16,\ \ x_L,x_P\ge0,
\]
cioè $\max\{c^Tx: Ax\le b,\ x\ge0\}$ con
$c=\binom{3}{5}$ e
\[
A=\begin{pmatrix}1&1\\1&0\\0&1\\1&2\end{pmatrix},\qquad b=\begin{pmatrix}12\\10\\6\\16\end{pmatrix}.
\]
(Dividere un vincolo per una costante positiva o riscalare l'obiettivo non cambia le soluzioni ottime; l'ottimo è ancora $(8,4)$, con valore $44$ migliaia di euro.)

\section{Definizioni (richiamo)}
Per $(P)\ \max\{c^Tx: Ax\le b, x\ge0\}$:
\begin{itemize}[nosep]
\item \textbf{soluzione ammissibile}: $\bar x$ con $A\bar x\le b$, $\bar x\ge0$;
\item \textbf{regione ammissibile}: $X=\{x\in\R^n: Ax\le b,\ x\ge0\}$;
\item \textbf{problema inammissibile}: $X=\emptyset$;
\item \textbf{soluzione ottima}: $x^*\in X$ con $c^Tx^*\ge c^Tx$ per ogni $x\in X$;
\item \textbf{problema illimitato}: per ogni $M$ esiste $x\in X$ con $c^Tx>M$.
\end{itemize}

\section{Forme di un problema di PL}

In generale un problema di PL può avere vincoli di tutti i tipi e variabili di tutti i segni:
\[
\max\ (\text{o}\min)\ c^Tx\quad\st\ A_1x\le b_1,\quad A_2x\ge b_2,\quad A_3x=b_3,\quad x_j\ge0,\ x_j\le0\ \text{o libere}.
\]
Due forme ``di riferimento'':
\begin{center}
\begin{tabular}{@{}c@{\qquad\qquad}c@{}}
\textbf{Forma standard} & \textbf{Forma canonica (o generale)}\\[3pt]
$\begin{aligned}\min\ &c^Tx\\ \st\ &Ax=b\\ &x\ge0\end{aligned}$ \quad (con $b\ge0$)
&
$\begin{aligned}\max\ &c^Tx\\ \st\ &Ax\le b\\ &x\ge0\end{aligned}$
\end{tabular}
\end{center}
La forma standard serve al \textbf{simplesso} (lavora con sistemi di equazioni), la forma canonica è comoda per la \textbf{geometria} e per la \textbf{dualità}.

\begin{prop}
Ogni problema di PL può essere scritto in modo equivalente in forma standard, e in modo equivalente in forma canonica.
\end{prop}
``Equivalente'' significa: da ogni soluzione (ottima) dell'uno si ricava in modo immediato una soluzione (ottima) dell'altro, con lo stesso valore dell'obiettivo (a meno del segno).

\subsection{Trasformazioni}
\begin{center}
\renewcommand{\arraystretch}{1.3}
\begin{tabular}{@{}p{0.2\textwidth}p{0.36\textwidth}p{0.36\textwidth}@{}}
\toprule
 & Da & A\\ \midrule
\textbf{Obiettivo} & $\max c^Tx$ & $-\min(-c^T x)$\\
 & $\min c^Tx$ & $-\max(-c^Tx)$\\ \midrule
\textbf{Variabili} & $x_j\le0$ & $x_j=-x_j'$, \ $x_j'\ge0$\\
 & $x_j$ libera & $x_j=x_j^+-x_j^-$, \ $x_j^+,x_j^-\ge0$\\ \midrule
\textbf{Vincoli} & $a^Tx\le b$ & $a^Tx+s=b$, \ $s\ge0$ \ (\emph{slack})\\
 & $a^Tx\ge b$ & $a^Tx-s=b$, \ $s\ge0$ \ (\emph{surplus})\\
 & $a^Tx\ge b$ & $-a^Tx\le -b$\\
 & $a^Tx=b$ & $a^Tx\le b$ \ e \ $-a^Tx\le -b$\\
 & $a^Tx=b$ con $b<0$ & $-a^Tx=-b$ \ (per avere $b\ge0$)\\
\bottomrule
\end{tabular}
\end{center}
Le variabili di slack/surplus hanno un significato: la slack del vincolo di terreno è l'\emph{ettaraggio non usato}. All'ottimo, slack $=0$ significa vincolo \emph{attivo} (risorsa satura).

\begin{esbox}[coltivatore in forma standard]
\[
\begin{aligned}
-\min\ & -3x_L-5x_P\\
\st\ & x_L+x_P+s_1=12\\
& x_L+s_2=10\\
& x_P+s_3=6\\
& x_L+2x_P+s_4=16\\
& x_L,x_P,s_1,s_2,s_3,s_4\ge0
\end{aligned}
\]
All'ottimo $(x_L,x_P)=(8,4)$: $s=(0,2,2,0)$ (terreno e stallatico saturi; avanzano 2 unità di semi e di tuberi nelle unità riscalate). Il coltivatore è \emph{già} in forma canonica.
\end{esbox}

\begin{esbox}[problema misto in entrambe le forme]
$\min\{x_1-2x_2:\ x_1+x_2\ge 2,\ \ x_1-x_2=-1,\ \ x_1\ge0,\ x_2\text{ libera}\}$.
\textbf{Standard}: $x_2=x_2^+-x_2^-$; il vincolo $\ge$ prende un surplus; l'uguaglianza ha termine noto negativo e si moltiplica per $-1$:
\[\min\ x_1-2x_2^++2x_2^-\ \ \st\ x_1+x_2^+-x_2^--s_1=2,\ \ -x_1+x_2^+-x_2^-=1,\ \ x_1,x_2^+,x_2^-,s_1\ge0.\]
\textbf{Canonica}: $\max\{-x_1+2x_2^+-2x_2^-:\ -x_1-x_2^++x_2^-\le-2,\ \ x_1-x_2^++x_2^-\le-1,\ \ -x_1+x_2^+-x_2^-\le1,\ \ x\ge0\}$ (il valore ottimo è l'opposto).
\end{esbox}

\section{Risoluzione grafica}

Ogni vincolo $a_i^Tx\le b_i$ è associato all'\textbf{iperpiano} $\{x\in\R^n:a_i^Tx=b_i\}$ (in $\R^2$ una retta) e definisce il \textbf{semispazio chiuso} ammissibile $\{x:a_i^Tx\le b_i\}$. La regione ammissibile è l'intersezione di tutti i semispazi.

\begin{center}
\begin{tikzpicture}[scale=0.5,>=Stealth]
\begin{scope}
\clip (-1,-1) rectangle (9,8.5);
\fill[blue!10] (-1,8.5)--(-1,-1)--(9,-1)--(9,3.5)--cycle;
\draw[thick] (-1,8.5)--(9,3.5);
\foreach \t in {0,0.08,...,1} \draw[gray!70] ($(-1,8.5)!\t!(9,3.5)$)--++(0.25,0.5);
\end{scope}
\draw[->] (-1,0)--(9.3,0) node[below] {$x_1$}; \draw[->] (0,-1)--(0,8.8) node[left] {$x_2$};
\draw[->,very thick,red!70!black] (4,5.5)--++(1,2) node[above,font=\scriptsize] {$a$};
\node[font=\scriptsize,fill=white,inner sep=1pt] at (3,2) {$a^Tx\le b$};
\node[font=\scriptsize,rotate=-26.6] at (6.2,5.5) {$a^Tx=b$};
\node[font=\scriptsize,align=center] at (4,-2.3) {(i) un vincolo: iperpiano e semispazio;\\ $a$ è normale all'iperpiano e punta fuori};
\begin{scope}[xshift=13cm]
\fill[blue!12] (0,0)--(5,0)--(5,1.5)--(3,3.5)--(0,4)--cycle;
\draw[gray] (-0.5,4.08)--(8,2.67); \draw[gray] (5,-0.6)--(5,7.5); \draw[gray] (-0.5,7.5)--(7.5,-0.5);
\draw[thick] (0,0)--(5,0)--(5,1.5)--(3,3.5)--(0,4)--cycle;
\draw[->] (-1,0)--(8.3,0) node[below] {$x_1$}; \draw[->] (0,-1)--(0,8.8) node[left] {$x_2$};
\node[font=\scriptsize,align=center] at (4,-2.3) {(ii) intersezione di semispazi\\ $=$ poligono convesso (poliedro)};
\end{scope}
\end{tikzpicture}
\end{center}
Per capire da che parte sta il semispazio basta provare un punto: se l'origine soddisfa il vincolo, la parte ammissibile è quella che contiene l'origine. Il vettore dei coefficienti $a$ punta sempre verso la parte \emph{non} ammissibile di un vincolo $\le$.

Le \textbf{curve di livello} dell'obiettivo sono gli iperpiani $c^Tx=z$: rette parallele, ognuna associata a un valore $z$. Il \textbf{gradiente} $\nabla(c^Tx)=c$ è perpendicolare alle curve di livello e indica la direzione di crescita.

\textbf{Procedura} (2 variabili): disegna la regione; disegna una curva di livello; spostala parallelamente nella direzione di $c$ (per un $\min$: nella direzione $-c$) fino all'ultima posizione in cui tocca ancora la regione. Il valore di quella retta è l'ottimo, i punti di contatto sono le soluzioni ottime. Per il coltivatore: $3x_L+5x_P=0,30,42,44$; l'ultima che tocca è $z=44$ in $(8,4)$. Il punto esatto si trova risolvendo il sistema dei due vincoli attivi: $x_L+x_P=12$, $x_L+2x_P=16$.

\begin{center}
\begin{tikzpicture}[scale=0.62,>=Stealth]
\begin{scope}
\clip (-0.8,-0.8) rectangle (14,9.6);
\fill[blue!12] (0,0)--(10,0)--(10,2)--(8,4)--(4,6)--(0,6)--cycle;
\draw[thick,red!70!black] (12,0)--(2.4,9.6);
\draw[thick,green!50!black] (10,-0.8)--(10,9.6);
\draw[thick,orange!80!black] (-0.8,6)--(14,6);
\draw[thick,violet] (-0.8,8.4)--(14,1);
\foreach \z/\c in {0/gray,30/gray,42/gray,44/black} \draw[dashed,\c] (\z/3+1.6,-0.96)--(-0.8,\z/5+0.48);
\end{scope}
\draw[->] (-0.8,0)--(14.3,0) node[below] {$x_L$}; \draw[->] (0,-0.8)--(0,9.9) node[left] {$x_P$};
\foreach \x in {2,4,...,12} \draw (\x,0.1)--(\x,-0.1) node[below,font=\scriptsize] {\x};
\foreach \y in {2,4,6,8} \draw (0.1,\y)--(-0.1,\y) node[left,font=\scriptsize] {\y};
\node[font=\scriptsize,gray,anchor=east] at (-0.9,0.2) {$z=0$};
\node[font=\scriptsize,gray,fill=blue!12,inner sep=1pt] at (5,3) {$z=30$};
\node[font=\scriptsize,gray,fill=white,inner sep=1pt] at (1.5,7.5) {$z=42$};
\node[font=\scriptsize,fill=white,inner sep=1pt] at (13.1,0.95) {$z=44$};
\draw[->,very thick] (1.5,1.5)--++(1.2,2) node[above,font=\scriptsize] {$c=(3,5)$};
\fill (8,4) circle (3.5pt);
\draw (8,4)--(11.6,5) node[right,font=\small,fill=white,inner sep=1pt] {$x^*=(8,4)$};
\fill (4,6) circle (2.5pt);
\begin{scope}[xshift=15.2cm,yshift=8.6cm,font=\scriptsize]
\draw[thick,red!70!black] (0,0)--(0.8,0) node[right,black] {$x_L+x_P\le12$ (terreno)};
\draw[thick,green!50!black] (0,-0.7)--(0.8,-0.7) node[right,black] {$x_L\le10$ (semi)};
\draw[thick,orange!80!black] (0,-1.4)--(0.8,-1.4) node[right,black] {$x_P\le6$ (tuberi)};
\draw[thick,violet] (0,-2.1)--(0.8,-2.1) node[right,black] {$x_L+2x_P\le16$ (stallatico)};
\draw[dashed] (0,-2.8)--(0.8,-2.8) node[right,black] {curve di livello $3x_L+5x_P=z$};
\end{scope}
\end{tikzpicture}
\end{center}
Le curve $z=0,30,42$ intersecano la regione; $z=44$ la tocca solo in $(8,4)$; ogni $z>44$ non la tocca più. Nota che il vertice $(4,6)$ (valore 42) non è ottimo anche se ``sta in alto'': conta la direzione di $c$, non la posizione.

\begin{esbox}[un problema di minimo con regione illimitata]
\[\min\ x_1+x_2\quad\st\ x_1+2x_2\ge4,\quad 3x_1+x_2\ge6,\quad x_1,x_2\ge0.\]
\begin{center}
\begin{tikzpicture}[scale=0.62,>=Stealth]
\begin{scope}
\clip (-0.6,-0.6) rectangle (7.5,7.5);
\fill[blue!12] (0,7.5)--(0,6)--(1.6,1.2)--(4,0)--(7.5,0)--(7.5,7.5)--cycle;
\draw[thick,red!70!black] (-0.6,7.8)--(2.5,-1.5);
\draw[thick,violet] (-0.6,2.3)--(7.5,-1.75);
\draw[dashed] (0,2.8)--(2.8,0);
\draw[dashed,gray] (0,5)--(5,0);
\draw[dashed,gray] (0,7)--(7,0);
\end{scope}
\draw[->] (-0.6,0)--(7.8,0) node[below] {$x_1$}; \draw[->] (0,-0.6)--(0,7.8) node[left] {$x_2$};
\foreach \x in {2,4,6} \draw (\x,0.1)--(\x,-0.1) node[below,font=\scriptsize] {\x};
\foreach \y in {2,4,6} \draw (0.1,\y)--(-0.1,\y) node[left,font=\scriptsize] {\y};
\fill (1.6,1.2) circle (3.5pt) node[left=3pt,font=\scriptsize,fill=white,inner sep=1pt] {$x^*=(1.6,\,1.2)$};
\fill (0,6) circle (2.5pt) node[right,font=\scriptsize] {$(0,6)$: 6};
\fill (4,0) circle (2.5pt) node[above right,font=\scriptsize] {$(4,0)$: 4};
\draw[->,very thick] (5,5)--(4,4) node[below left,font=\scriptsize] {$-c$};
\node[font=\scriptsize,red!70!black] at (2.6,6.8) {$3x_1+x_2=6$};
\node[font=\scriptsize,violet,anchor=east] at (-0.7,1.7) {$x_1+2x_2=4$};
\node[font=\scriptsize,gray] at (5.6,2.1) {$z=5$};
\end{tikzpicture}
\end{center}
Per un minimo ci si muove lungo $-c=(-1,-1)$: le curve di livello $z=7,5,\dots$ scendono finché toccano la regione per l'ultima volta nel vertice $(1.6,1.2)$, intersezione dei due vincoli ($x_1+2x_2=4$, $3x_1+x_2=6$), con $z^*=2.8$. Confronto con gli altri vertici: $(0,6)$ vale 6, $(4,0)$ vale 4.
\textbf{Regione illimitata ma ottimo finito}: le direzioni di recessione sono $d\ge0$ e $c^Td=d_1+d_2\ge0$, quindi muoversi all'infinito fa solo \emph{peggiorare} un minimo. Lo stesso problema con $\max$ sarebbe invece illimitato.
\end{esbox}

\textbf{Osservazione}: la soluzione ottima si trova sulla \emph{frontiera} della regione ammissibile.

\subsection{Quante soluzioni ottime può avere una PL?}

\begin{center}
\begin{tikzpicture}[scale=0.36,>=Stealth]
\begin{scope}
\fill[blue!12] (0,0)--(10,0)--(10,2)--(8,4)--(4,6)--(0,6)--cycle;
\draw[->] (-0.3,0)--(11.5,0); \draw[->] (0,-0.3)--(0,8);
\draw[thick] (0,0)--(10,0)--(10,2)--(8,4)--(4,6)--(0,6)--cycle;
\draw[dashed,red] (3,7)--(11.5,1.9); \fill[red] (8,4) circle (5pt);
\node[font=\scriptsize,align=center] at (5.5,-2) {(a) ottimo unico\\ $\max 3x_1+5x_2$};
\end{scope}
\begin{scope}[xshift=14cm]
\fill[blue!12] (0,0)--(10,0)--(10,2)--(8,4)--(4,6)--(0,6)--cycle;
\draw[->] (-0.3,0)--(11.5,0); \draw[->] (0,-0.3)--(0,8);
\draw[thick] (0,0)--(10,0)--(10,2)--(8,4)--(4,6)--(0,6)--cycle;
\draw[ultra thick,red] (8,4)--(10,2);
\node[font=\scriptsize,align=center] at (5.5,-2) {(b) infinite soluzioni ottime\\ $\max x_1+x_2$ (spigolo)};
\end{scope}
\begin{scope}[xshift=28cm]
\fill[blue!12] (0,0)--(11,0)--(11,5)--(0,5)--cycle;
\draw[->] (-0.3,0)--(11.5,0); \draw[->] (0,-0.3)--(0,8);
\draw[thick] (11,0)--(0,0)--(0,5)--(11,5);
\draw[->,thick,red] (4,2.5)--(8,2.5);
\node[font=\scriptsize,align=center] at (5.5,-2) {(c) illimitato\\ $\max x_1$, \ $0\le x_2\le5$, $x_1\ge0$};
\end{scope}
\begin{scope}[yshift=-12.5cm]
\fill[blue!12] (0,0)--(10,0)--(10,2)--(8,4)--(4,6)--(0,6)--cycle;
\draw[->] (-0.3,0)--(11.5,0); \draw[->] (0,-0.3)--(0,8);
\draw[thick] (0,0)--(10,0)--(10,2)--(8,4)--(4,6)--(0,6)--cycle;
\draw[dashed,red] (-1,0.6)--(2,-1.2);
\fill[red] (0,0) circle (5pt);
\draw[->,thick,red] (6,4)--(4.8,2);
\node[font=\scriptsize,align=center] at (5.5,-2) {(a') ottimo unico (minimo)\\ $\min 3x_1+5x_2$: $x^*=(0,0)$};
\end{scope}
\begin{scope}[xshift=14cm,yshift=-12.5cm]
\fill[blue!12] (0,8)--(0,3)--(3,1)--(8,0)--(11,0)--(11,8)--cycle;
\draw[->] (-0.3,0)--(11.5,0); \draw[->] (0,-0.3)--(0,8);
\draw[thick] (0,8)--(0,3)--(3,1)--(8,0)--(11,0);
\draw[ultra thick,red] (0,8)--(0,3);
\draw[->,thick,red] (4,4)--(2.5,4);
\node[font=\scriptsize,align=center] at (5.5,-2) {(b') infiniti ottimi (semiretta)\\ $\min x_1$: lato $x_1=0$, $x_2\ge3$};
\end{scope}
\begin{scope}[xshift=28cm,yshift=-12.5cm]
\fill[blue!20] (0,0)--(2.5,0)--(0,2.5)--cycle;
\fill[red!15] (5,0)--(11,0)--(11,8)--(0,8)--(0,5)--cycle;
\draw[thick,blue!60!black] (-0.3,2.8)--(2.8,-0.3);
\draw[thick,red!70!black] (-0.3,5.3)--(5.3,-0.3);
\draw[->] (-0.3,0)--(11.5,0); \draw[->] (0,-0.3)--(0,8);
\node[font=\scriptsize,blue!60!black] at (1.2,0.7) {$\le1$};
\node[font=\scriptsize,red!70!black] at (5.5,4) {$x_1+x_2\ge2$};
\node[font=\scriptsize,align=center] at (5.5,-2) {(d) inammissibile\\ i due semispazi non si intersecano};
\end{scope}
\end{tikzpicture}
\end{center}
\begin{itemize}[nosep]
\item \textbf{(a), (a') Unica} soluzione ottima, in un vertice. Con lo stesso $c$, il massimo e il minimo stanno in vertici ``opposti'' rispetto alla direzione $c$.
\item \textbf{(b) Infinite} soluzioni ottime: la curva di livello è parallela a un lato della regione; ottimi tutti i punti del segmento tra $(8,4)$ e $(10,2)$, valore 12. Gli estremi del segmento sono vertici. In (b') l'insieme degli ottimi è addirittura illimitato (una semiretta), ma il valore ottimo ($0$) è finito.
\item \textbf{(c) Problema illimitato}: regione illimitata e obiettivo che cresce lungo una direzione in cui ci si può muovere all'infinito. Attenzione: \emph{regione illimitata non implica problema illimitato}: sulla stessa regione $\min x_1$ ha ottimo (segmento $x_1=0$) e $\max x_2$ ha ottimo $x_2=5$ (infiniti punti).
\item \textbf{Problema inammissibile}: $X=\emptyset$, nessuna soluzione (es.\ $x_1+x_2\le1$, $x_1+x_2\ge2$, $x\ge0$: figura sotto, (d)).
\end{itemize}
\begin{esame}
Un problema di PL è sempre in uno di questi casi: \textbf{inammissibile}, \textbf{illimitato}, oppure ha \textbf{ottimo finito}, raggiunto in 1 punto o in infiniti punti. \textbf{Mai esattamente 2} (o un numero finito $>1$) soluzioni ottime. Con variabili intere invece può succedere: $\max\{x_1+x_2: x_1+x_2\le1,\ x\in\{0,1\}^2\}$ ha esattamente 2 soluzioni ottime.
\end{esame}

\section{Geometria della PL}

\subsection{Combinazioni}
Dati $x^1,\dots,x^k\in\R^n$ e pesi $\lambda_1,\dots,\lambda_k$, il punto $y=\sum_i\lambda_ix^i$ è una combinazione
\begin{center}
\begin{tabular}{@{}lll@{}}
\toprule
\textbf{lineare} & $\lambda_i\in\R$ & in $\R^2$, di 2 vettori indipendenti: tutto il piano\\
\textbf{affine} & $\sum_i\lambda_i=1$ & di 2 punti: la retta che li contiene\\
\textbf{conica} & $\lambda_i\ge0$ & di 2 vettori: l'``angolo'' tra le due semirette dall'origine\\
\textbf{convessa} & $\lambda_i\ge0$, $\sum_i\lambda_i=1$ & di 2 punti: il segmento; di 3: il triangolo\\
\bottomrule
\end{tabular}
\end{center}

\begin{esbox}[esempi delle slide]
\begin{itemize}[nosep]
\item $(2,2)$ è combinazione convessa di $(4,0)$ e $(1,3)$: $\lambda(4,0)+(1-\lambda)(1,3)=(1+3\lambda,\ 3-3\lambda)=(2,2)$ per $\lambda=\tfrac13\in[0,1]$.
\item $(2,2)$ è combinazione convessa di $(1,1),(3,1),(2,3)$: dal sistema $\lambda_1+3\lambda_2+2\lambda_3=2$, $\lambda_1+\lambda_2+3\lambda_3=2$, $\lambda_1+\lambda_2+\lambda_3=1$ si ottiene $\lambda=(\tfrac14,\tfrac14,\tfrac12)\ge0$.
\item $(4,4)$ è combinazione conica di $(2,3)$ e $(3,1)$: $2a+3b=4$, $3a+b=4$ danno $a=\tfrac87$, $b=\tfrac47$, entrambi $\ge0$. (Non è convessa: $a+b=\tfrac{12}{7}\ne1$.)
\end{itemize}
Metodo: imposta il sistema lineare nei pesi e controlla i segni (e la somma).
\end{esbox}

\textbf{Le quattro combinazioni di due vettori} $x'=(4,0)$ e $x''=(1,3)$ (in rosso l'insieme di tutte le combinazioni del tipo indicato):
\begin{center}
\begin{tikzpicture}[scale=0.45,>=Stealth]
\foreach \k/\name in {0/lineare,1/affine,2/conica,3/convessa}{
\begin{scope}[xshift=\k*8.2cm]
\ifcase\k
  \fill[red!15] (-1.5,-1.5) rectangle (5.5,5.5);
\or
  \draw[very thick,red!70!black] (5.33,-1)--(-0.33,4.67);
\or
  \fill[red!15] (0,0)--(5.5,0)--(5.5,5.5)--(1.83,5.5)--cycle;
  \draw[very thick,red!70!black] (5.5,0)--(0,0)--(1.83,5.5);
\or
  \draw[very thick,red!70!black] (4,0)--(1,3);
\fi
\draw[->,gray] (-1.5,0)--(5.7,0); \draw[->,gray] (0,-1.5)--(0,5.7);
\draw[->,thick] (0,0)--(4,0); \draw[->,thick] (0,0)--(1,3);
\fill (4,0) circle (3pt) node[below,font=\scriptsize] {$x'$};
\fill (1,3) circle (3pt) node[left,font=\scriptsize] {$x''$};
\node[font=\small] at (2,-2.6) {\name};
\end{scope}}
\end{tikzpicture}
\end{center}
Lineare: tutto il piano (i due vettori sono indipendenti). Affine: la retta per i due punti (anche oltre gli estremi, con un peso negativo). Conica: l'angolo tra le due semirette uscenti dall'origine. Convessa: solo il segmento. Il punto $(2,2)$ sta sul segmento, quindi è combinazione convessa (e anche affine, conica e lineare).

\begin{center}
\begin{tikzpicture}[scale=0.62,>=Stealth]
\begin{scope}
\fill[blue!12] (1,1)--(3,1)--(2,3)--cycle; \draw[thick] (1,1)--(3,1)--(2,3)--cycle;
\foreach \p/\l/\a in {{(1,1)}/{(1,1)}/below left,{(3,1)}/{(3,1)}/below right,{(2,3)}/{(2,3)}/above} \fill \p circle (2.5pt) node[\a,font=\scriptsize] {$\l$};
\fill[red] (2,2) circle (3pt) node[right=2pt,font=\scriptsize,red] {$(2,2)=\tfrac14(1,1)+\tfrac14(3,1)+\tfrac12(2,3)$};
\draw[->,gray] (-0.3,0)--(4,0); \draw[->,gray] (0,-0.3)--(0,4);
\node[font=\scriptsize] at (2,-0.9) {combinazione convessa di 3 punti: il triangolo};
\end{scope}
\begin{scope}[xshift=10cm]
\fill[blue!12] (0,0)--(5.5,1.83)--(5.5,5.5)--(3.67,5.5)--cycle;
\draw[->,thick] (0,0)--(2,3) node[left,font=\scriptsize] {$(2,3)$};
\draw[->,thick] (0,0)--(3,1) node[below right,font=\scriptsize] {$(3,1)$};
\draw[dashed] (2.286,3.429)--(4,4)--(1.714,0.571);
\draw[->,red!70!black] (0,0)--(2.286,3.429) node[left,font=\scriptsize] {$\tfrac87(2,3)$};
\draw[->,red!70!black] (0,0)--(1.714,0.571);
\node[font=\scriptsize,red!70!black] at (2.3,0.05) {$\tfrac47(3,1)$};
\fill[red] (4,4) circle (3pt) node[right,font=\scriptsize,red] {$(4,4)$};
\draw[->,gray] (-0.3,0)--(5.7,0); \draw[->,gray] (0,-0.3)--(0,5.7);
\node[font=\scriptsize] at (2.8,-0.9) {combinazione conica: regola del parallelogramma};
\end{scope}
\end{tikzpicture}
\end{center}

\subsection{Insiemi convessi, coni, involucri}
\begin{defbox}[insieme convesso]
$S\subseteq\R^n$ è \textbf{convesso} se per ogni $x,y\in S$ e $\lambda\in[0,1]$ si ha $\lambda x+(1-\lambda)y\in S$ (contiene il segmento tra due suoi punti qualsiasi).
\end{defbox}
\begin{defbox}[cono]
$C\subseteq\R^n$ è un \textbf{cono} se $x\in C\Rightarrow \lambda x\in C$ per ogni $\lambda\ge0$. Un cono non è necessariamente convesso: l'unione dei due semiassi positivi $\{x\ge0: x_1x_2=0\}$ è un cono non convesso; il primo ortante $\R^2_+$ è un cono convesso.
\end{defbox}
\begin{defbox}[involucro convesso e conico]
L'\textbf{involucro (inviluppo) convesso} $\mathrm{conv}(K)$ è l'insieme di tutte le combinazioni convesse di elementi di $K$: il più piccolo convesso che contiene $K$. L'\textbf{involucro conico} $\mathrm{cono}(K)$ è l'insieme di tutte le combinazioni coniche: il più piccolo cono convesso che contiene $K$.
\end{defbox}
Esempio: $\mathrm{conv}\{(0,0),(2,0),(1,2)\}$ è il triangolo pieno con quei vertici, cioè $\{x: x_2\ge0,\ x_2\le2x_1,\ x_2\le4-2x_1\}$.

\needspace{12\baselineskip}
\paragraph{Insiemi convessi e non convessi.}
\begin{center}
\begin{tikzpicture}[scale=0.55,>=Stealth]
\begin{scope}
\fill[blue!12] (2,1.5) ellipse (2 and 1.4); \draw[thick] (2,1.5) ellipse (2 and 1.4);
\draw[red!70!black,thick] (0.7,1.0)--(3.3,2.2); \fill (0.7,1) circle (2.5pt) (3.3,2.2) circle (2.5pt);
\node[font=\scriptsize] at (2,-0.6) {convesso};
\end{scope}
\begin{scope}[xshift=5.5cm]
\fill[blue!12] (0,0)--(4,0)--(2.5,3)--cycle; \draw[thick] (0,0)--(4,0)--(2.5,3)--cycle;
\draw[red!70!black,thick] (1,0.3)--(2.6,2.4); \fill (1,0.3) circle (2.5pt) (2.6,2.4) circle (2.5pt);
\node[font=\scriptsize] at (2,-0.6) {convesso};
\end{scope}
\begin{scope}[xshift=11.5cm]
\fill[blue!12] (0,0)--(4,0)--(4,1)--(1,1)--(1,3)--(0,3)--cycle; \draw[thick] (0,0)--(4,0)--(4,1)--(1,1)--(1,3)--(0,3)--cycle;
\draw[red!70!black,thick,dashed] (0.5,2.6)--(3.6,0.6); \fill (0.5,2.6) circle (2.5pt) (3.6,0.6) circle (2.5pt);
\node[font=\scriptsize,align=center] at (2,-0.9) {non convesso:\\il segmento esce};
\end{scope}
\begin{scope}[xshift=17.5cm]
\fill[blue!12] (2,1.5) circle (1.5); \fill[white] (2,1.5) circle (0.7); \draw[thick] (2,1.5) circle (1.5) (2,1.5) circle (0.7);
\draw[red!70!black,thick,dashed] (0.8,1.2)--(3.2,1.8); \fill (0.8,1.2) circle (2.5pt) (3.2,1.8) circle (2.5pt);
\node[font=\scriptsize,align=center] at (2,-0.9) {non convesso\\(corona)};
\end{scope}
\end{tikzpicture}
\end{center}
La corona circolare $\{1\le\|x\|\le2\}$ è un esempio di insieme definito da un vincolo ``$\ge$'' su una funzione convessa: non è convesso. Un vincolo $\|x\|\le2$ invece dà un insieme convesso.

\needspace{12\baselineskip}
\paragraph{Coni.}
\begin{center}
\begin{tikzpicture}[scale=0.52,>=Stealth]
\begin{scope}
\fill[blue!12] (0,0)--(4,1.2)--(4,4)--(1.6,4)--cycle;
\draw[thick] (4,1.2)--(0,0)--(1.6,4);
\fill (0,0) circle (2.5pt) node[below left,font=\scriptsize] {$0$};
\fill[red] (1.5,1.3) circle (2.5pt); \fill[red] (3,2.6) circle (2.5pt); \draw[->,red,dashed] (1.5,1.3)--(2.9,2.51);
\node[font=\scriptsize,red] at (2.2,2.4) {$\lambda x$};
\node[font=\scriptsize,align=center] at (2,-1) {cono convesso\\(poliedrico)};
\end{scope}
\begin{scope}[xshift=6cm]
\draw[very thick,blue!60!black] (0,0)--(4,0); \draw[very thick,blue!60!black] (0,0)--(0,4);
\draw[red!70!black,dashed,thick] (2.5,0)--(0,2.5); \fill[red!70!black] (1.25,1.25) circle (2.5pt);
\fill (0,0) circle (2.5pt);
\node[font=\scriptsize,align=center] at (2,-1) {cono non convesso\\(due semiassi)};
\end{scope}
\begin{scope}[xshift=12cm]
\fill[blue!12] (0,0)--(4,0.8)--(4,2.4)--cycle; \fill[blue!12] (0,0)--(1.2,4)--(-1.2,4)--cycle;
\draw[thick] (4,0.8)--(0,0)--(4,2.4) (1.2,4)--(0,0)--(-1.2,4);
\fill (0,0) circle (2.5pt);
\node[font=\scriptsize,align=center] at (1.4,-1) {cono non convesso\\(due settori)};
\end{scope}
\begin{scope}[xshift=18cm]
\fill[blue!12] (0,0)--(3,1)--(1,3)--cycle; \draw[thick] (0,0)--(3,1)--(1,3)--cycle;
\fill (1.2,1.2) circle (2.5pt) node[right,font=\scriptsize] {$x$};
\fill[red] (3.6,3.6) circle (2.5pt) node[right,font=\scriptsize,red] {$3x\notin$};
\fill (0,0) circle (2.5pt);
\node[font=\scriptsize,align=center] at (1.8,-1) {\emph{non} è un cono\\(limitato)};
\end{scope}
\begin{scope}[xshift=26cm]
\fill[blue!12] (0,0)--(-1.8,3.6)--(1.8,3.6)--cycle;
\fill[blue!20] (0,3.6) ellipse (1.8 and 0.45); \draw[thick] (0,3.6) ellipse (1.8 and 0.45);
\draw[thick] (-1.8,3.6)--(0,0)--(1.8,3.6);
\fill (0,0) circle (2.5pt);
\node[font=\scriptsize,align=center] at (0,-1) {cono di secondo ordine\\ in $\R^3$: convesso,\\ non poliedrico};
\end{scope}
\end{tikzpicture}
\end{center}
Un cono contiene sempre l'origine (con $\lambda=0$) ed è ``illimitato'' in ogni sua direzione: l'unico cono limitato è $\{0\}$. Un cono convesso in $\R^2$ diverso dal piano e dal semipiano è sempre un ``angolo'' (settore) con vertice nell'origine.

\needspace{12\baselineskip}
\paragraph{Involucri.}
\begin{center}
\begin{tikzpicture}[scale=0.6,>=Stealth]
\begin{scope}
\fill[blue!12] (0,0)--(2,0)--(1,2)--cycle; \draw[thick] (0,0)--(2,0)--(1,2)--cycle;
\foreach \p in {(0,0),(2,0),(1,2)} \fill \p circle (2.5pt);
\node[font=\scriptsize,align=center] at (1,-1) {$\mathrm{conv}\{(0,0),$\\$(2,0),(1,2)\}$};
\end{scope}
\begin{scope}[xshift=4.6cm]
\fill[blue!12] (0,0.5)--(1.5,0)--(3.5,0.6)--(4,2.2)--(2.2,3.2)--(0.4,2.4)--cycle;
\draw[thick] (0,0.5)--(1.5,0)--(3.5,0.6)--(4,2.2)--(2.2,3.2)--(0.4,2.4)--cycle;
\foreach \p in {(0,0.5),(1.5,0),(3.5,0.6),(4,2.2),(2.2,3.2),(0.4,2.4)} \fill \p circle (2.5pt);
\foreach \p in {(1.2,1.2),(2.5,1.8),(2,0.9),(3,1.5),(1.3,2.3)} \fill[gray] \p circle (2pt);
\node[font=\scriptsize,align=center] at (2,-0.8) {$\mathrm{conv}$ di un insieme finito:\\ gli ``esterni'' sono i vertici};
\end{scope}
\begin{scope}[xshift=10.6cm]
\fill[blue!12] (0,0)--(4,2)--(4,4)--(2,4)--cycle;
\draw[thick] (4,2)--(0,0)--(2,4);
\draw[->,very thick] (0,0)--(2,1) node[below right,font=\scriptsize] {$(2,1)$};
\draw[->,very thick] (0,0)--(1,2) node[left,font=\scriptsize] {$(1,2)$};
\node[font=\scriptsize] at (2,-1) {$\mathrm{cono}\{(2,1),(1,2)\}$};
\end{scope}
\begin{scope}[xshift=16.8cm]
\fill[blue!12] (0,0)--(4,0)--(4,4)--(0,4)--cycle;
\draw[->,very thick] (0,0)--(2,0) node[above right,font=\scriptsize] {$(1,0)$};
\draw[->,very thick] (0,0)--(0,2) node[left,font=\scriptsize] {$(0,1)$};
\draw[->,very thick] (0,0)--(1.4,1.4) node[above right,font=\scriptsize] {$(1,1)$};
\node[font=\scriptsize,align=center] at (2,-1.2) {$\mathrm{cono}\{(1,0),(0,1),(1,1)\}$\\ $=\R^2_+$ ($(1,1)$ superfluo)};
\end{scope}
\end{tikzpicture}
\end{center}

\subsection{Poliedri}
\textbf{Proprietà.}
\begin{itemize}[nosep]
\item L'intersezione di insiemi convessi (anche infiniti) è convessa.
\item Per $f$ lineare, $\{x: f(x)\le b\}$ è convesso per ogni $b$; $\{x\in\R^n: a^Tx\le b\}$ ($a\ne0$) è un \textbf{semispazio chiuso}.
\end{itemize}
\begin{defbox}[poliedro e politopo]
Un \textbf{poliedro} in $\R^n$ è l'intersezione di un numero \emph{finito} di semispazi chiusi; equivalentemente l'insieme delle soluzioni di un sistema di $m$ disequazioni lineari, $P=\{x\in\R^n: Ax\le b\}$. Un \textbf{politopo} è un poliedro limitato.
\end{defbox}
Quindi \textbf{la regione ammissibile di una PL è un poliedro} (convesso e chiuso); i vincoli $x\ge0$ sono semispazi come gli altri ($-x_j\le0$), le uguaglianze sono coppie di semispazi.

\emph{E con un vincolo non lineare?} In generale no: $\{x: x_1^2+x_2^2\le1\}$ è convesso ma non è un poliedro (servirebbero infiniti semispazi); $\{x: x_1^2+x_2^2\ge1\}$ non è nemmeno convesso.

\begin{center}
\begin{tikzpicture}[scale=0.55,>=Stealth]
\begin{scope}
\fill[blue!12] (0,0)--(3,0)--(4,1.5)--(2.5,3.2)--(0.3,2.5)--cycle; \draw[thick] (0,0)--(3,0)--(4,1.5)--(2.5,3.2)--(0.3,2.5)--cycle;
\node[font=\scriptsize,align=center] at (2,-1) {politopo\\ (poliedro limitato)};
\end{scope}
\begin{scope}[xshift=6.5cm]
\fill[blue!12] (0,4)--(0,1.5)--(1.5,0)--(4.5,0)--(4.5,4)--cycle; \draw[thick] (0,4)--(0,1.5)--(1.5,0)--(4.5,0);
\draw[->,red!70!black] (2.5,2)--(4,3.2);
\node[font=\scriptsize,align=center] at (2.2,-1) {poliedro illimitato};
\end{scope}
\begin{scope}[xshift=13.5cm]
\fill[blue!12] (2,1.6) circle (1.6); \draw[thick] (2,1.6) circle (1.6);
\draw[gray,dashed] (3.6,1.6)--(2.8,2.99)--(1.2,2.99)--(0.4,1.6)--(1.2,0.21)--(2.8,0.21)--cycle;
\node[font=\scriptsize,align=center] at (2,-1) {disco: convesso,\\ non poliedro};
\end{scope}
\begin{scope}[xshift=19.5cm]
\fill[blue!12] (0,0) rectangle (4,3.2); \fill[white] (2,1.6) circle (1); \draw[thick] (0,0) rectangle (4,3.2); \draw[thick] (2,1.6) circle (1);
\node[font=\scriptsize,align=center] at (2,-1) {$\{x_1^2+x_2^2\ge1\}\cap$ box:\\ non convesso};
\end{scope}
\end{tikzpicture}
\end{center}
Il disco si può approssimare con poligoni inscritti (tratteggio) ma nessun numero finito di semispazi lo descrive esattamente.

\textbf{Coni e poliedri.} Non tutti i coni convessi sono poliedri: il cono di secondo ordine (``gelato'') $\{x\in\R^3: x_3\ge\sqrt{x_1^2+x_2^2}\}$ è un cono convesso ma non poliedrico.
\begin{prop}
$P$ è un \textbf{cono poliedrico} se e solo se esiste una matrice $Q$ tale che $P=\{x\in\R^n: Qx\le0\}$ (termini noti tutti nulli).
\end{prop}

\subsection{Vertici e direzioni di recessione}
\begin{defbox}[vertice]
$v\in P$ è un \textbf{vertice} (punto estremo) di $P$ se non si può scrivere come combinazione convessa \emph{stretta} di due punti distinti di $P$: se $v=\lambda x+(1-\lambda)y$ con $x,y\in P$ e $\lambda\in(0,1)$, allora $x=y=v$.
\end{defbox}
Caratterizzazione algebrica (utile per il simplesso): in $P=\{x\in\R^n:Ax\le b\}$, $v$ è un vertice se e solo se tra i vincoli \emph{attivi} in $v$ ($a_i^Tv=b_i$) ce ne sono $n$ linearmente indipendenti. In $\R^2$: un vertice è l'intersezione di due rette di vincolo non parallele. Il coltivatore ha 6 vertici: $(0,0),(10,0),(10,2),(8,4),(4,6),(0,6)$. Il punto $(12,0)$ è intersezione di due rette di vincolo ma non è ammissibile, quindi \emph{non} è un vertice.

\begin{center}
\begin{tikzpicture}[scale=0.45,>=Stealth]
\fill[blue!12] (0,0)--(10,0)--(10,2)--(8,4)--(4,6)--(0,6)--cycle;
\draw[gray] (12,0)--(3,9); \draw[gray] (10,-0.6)--(10,8.5); \draw[gray] (-0.6,6)--(13,6); \draw[gray] (-0.6,8.3)--(13,1.5);
\draw[thick] (0,0)--(10,0)--(10,2)--(8,4)--(4,6)--(0,6)--cycle;
\draw[->] (-0.6,0)--(13.5,0) node[below] {$x_1$}; \draw[->] (0,-0.6)--(0,9) node[left] {$x_2$};
\foreach \p in {(0,0),(10,0),(10,2),(8,4),(4,6),(0,6)} \fill[red!70!black] \p circle (5pt);
\foreach \p in {(12,0),(6,6),(10,3),(10,6)} \draw[thick] \p ++(-0.25,-0.25)--++(0.5,0.5) \p ++(-0.25,0.25)--++(0.5,-0.5);
\fill (5,3) circle (3pt) node[right,font=\scriptsize] {interno};
\fill (9,3) circle (3pt);
\draw[dashed] (8,4)--(10,2); \node[font=\scriptsize,align=left,anchor=west] at (13.3,3.3) {$\bullet$ rossi: vertici (2 vincoli attivi indipendenti)\\ $\times$: intersezioni di rette di vincolo\\ \quad non ammissibili $\Rightarrow$ non vertici\\ $(9,3)$: sul bordo ma non vertice,\\ \quad è il punto medio di $(8,4)$ e $(10,2)$};
\end{tikzpicture}
\end{center}

\begin{defbox}[direzione di recessione]
$d\ne0$ è una \textbf{direzione di recessione} di $P$ se per ogni $x\in P$ e ogni $\lambda\ge0$ si ha $x+\lambda d\in P$ (da qualunque punto ci si può muovere all'infinito lungo $d$ restando in $P$). Per $P=\{x:Ax\le b\}\neq\emptyset$: $d$ è di recessione se e solo se $Ad\le0$. L'insieme $\{d: Ad\le0\}$ è il \emph{cono di recessione}.
\end{defbox}
Un politopo non ha direzioni di recessione. Esempio delle slide: $P=\{x\ge0,\ x_2\le5\}$: $Ad\le0$ dà $-d_1\le0$, $-d_2\le0$, $d_2\le0$, cioè $d=(d_1,0)$ con $d_1>0$: l'unica direzione (a meno di scala) è $(1,0)$.

\begin{center}
\begin{tikzpicture}[scale=0.5,>=Stealth]
\begin{scope}
\fill[blue!12] (0,0) rectangle (9,5); \draw[thick] (9,0)--(0,0)--(0,5)--(9,5);
\draw[->] (-0.5,0)--(9.5,0) node[below] {$x_1$}; \draw[->] (0,-0.5)--(0,6.5) node[left] {$x_2$};
\foreach \p in {(1,1),(2,4),(4,2.5)} {\fill \p circle (3pt); \draw[->,thick,red!70!black] \p--++(2.5,0);}
\draw[->,thick,gray,dashed] (5,3.5)--++(1.5,1.5); \node[font=\scriptsize,gray] at (7.3,5.6) {$(1,1)$: esce};
\node[font=\scriptsize] at (4.5,-1.3) {$\{x\ge0,\ x_2\le5\}$: cono di recessione $=\mathrm{cono}\{(1,0)\}$};
\end{scope}
\begin{scope}[xshift=13cm]
\fill[blue!12] (0,6.5)--(0,3)--(2,1)--(5,0)--(9,0)--(9,6.5)--cycle; \draw[thick] (0,6.5)--(0,3)--(2,1)--(5,0)--(9,0);
\draw[->] (-0.5,0)--(9.5,0) node[below] {$x_1$}; \draw[->] (0,-0.5)--(0,7) node[left] {$x_2$};
\begin{scope}[shift={(4.5,3)}]
\fill[red!15] (0,0)--(2.5,0)--(2.5,2.5)--(0,2.5)--cycle;
\draw[->,thick,red!70!black] (0,0)--(2.5,0); \draw[->,thick,red!70!black] (0,0)--(0,2.5);
\end{scope}
\node[font=\scriptsize] at (4.5,-1.3) {poliedro con $x\ge0$: cono di recessione $=\R^2_+$};
\end{scope}
\end{tikzpicture}
\end{center}
Il cono di recessione si può disegnare ``appoggiato'' in qualunque punto del poliedro: da ogni punto ci si muove lungo tutte le sue direzioni senza uscire.

\subsection{Teorema di decomposizione dei poliedri}
\begin{prop}[decomposizione, Minkowski--Weyl]
Ogni poliedro $P$ si può scrivere come
\[ P=\mathrm{conv}\{v^1,\dots,v^m\}+\mathrm{cono}\{e^1,\dots,e^p\}, \]
cioè ogni $x\in P$ è $x=\sum_{i}\lambda_iv^i+\sum_j\mu_je^j$ con $\lambda_i\ge0$, $\sum_i\lambda_i=1$, $\mu_j\ge0$: somma di un \textbf{politopo} e di un \textbf{cono} generato dalle direzioni di recessione. Se $P$ ha vertici, si possono prendere come $v^i$ esattamente i vertici. Vale anche il viceversa: un insieme di questa forma è un poliedro.
\end{prop}
(La ``somma'' è di Minkowski: $A+B=\{a+b: a\in A, b\in B\}$.) Il politopo o il cono possono essere banali: se $P$ è un politopo, $P=\mathrm{conv}(V)$ con $V$ insieme dei vertici.

\begin{esbox}[decomposizioni delle slide]
\begin{itemize}[nosep]
\item $P_1=\{x\in\R^2: x_1\ge1,\ x_2\ge1,\ x_1+x_2\ge3\}=\mathrm{conv}\{(1,2),(2,1)\}+\mathrm{cono}\{(1,0),(0,1)\}$.
\item $P_2=\{1\le x_1\le4,\ 1\le x_2\le3\}=\mathrm{conv}\{(1,1),(4,1),(4,3),(1,3)\}$ (politopo, cono $=\{0\}$).
\item $P_3=\{x_1\ge0,\ 0\le x_2\le1\}=\mathrm{conv}\{(0,0),(0,1)\}+\mathrm{cono}\{(1,0)\}$.
\item $P_4=\{0\le x_2\le1\}$ (striscia): \emph{nessun vertice} (contiene rette). $P_4=\mathrm{conv}\{(0,0),(0,1)\}+\mathrm{cono}\{(1,0),(-1,0)\}$: i punti $v^i$ esistono ma non sono vertici.
\end{itemize}
\end{esbox}
\begin{center}
\begin{tikzpicture}[scale=0.62,>=Stealth]
\begin{scope}
\fill[blue!12] (1,4.5)--(1,2)--(2,1)--(4.5,1)--(4.5,4.5)--cycle;
\draw[thick] (1,4.5)--(1,2)--(2,1)--(4.5,1);
\draw[ultra thick,red!70!black] (1,2)--(2,1);
\draw[->,very thick,green!50!black] (2,1)--(3.4,1); \draw[->,very thick,green!50!black] (1,2)--(1,3.4);
\draw[->] (-0.3,0)--(4.8,0); \draw[->] (0,-0.3)--(0,4.8);
\foreach \x in {1,2,3,4} \draw (\x,0.08)--(\x,-0.08) node[below,font=\tiny] {\x};
\foreach \y in {1,2,3,4} \draw (0.08,\y)--(-0.08,\y) node[left,font=\tiny] {\y};
\node[font=\small] at (3,3) {$P_1$};
\end{scope}
\begin{scope}[xshift=6.5cm]
\fill[blue!12] (1,1) rectangle (4,3); \draw[thick] (1,1) rectangle (4,3);
\foreach \p in {(1,1),(4,1),(4,3),(1,3)} \fill[red!70!black] \p circle (3pt);
\draw[->] (-0.3,0)--(4.8,0); \draw[->] (0,-0.3)--(0,4.8);
\node[font=\small] at (2.5,2) {$P_2$};
\end{scope}
\begin{scope}[xshift=13cm]
\fill[blue!12] (0,0) rectangle (4.5,1); \draw[thick] (4.5,0)--(0,0)--(0,1)--(4.5,1);
\draw[ultra thick,red!70!black] (0,0)--(0,1);
\draw[->,very thick,green!50!black] (0,0.5)--(1.6,0.5);
\draw[->] (-0.3,0)--(4.8,0); \draw[->] (0,-0.3)--(0,4.8);
\node[font=\small] at (3,2) {$P_3$};
\end{scope}
\begin{scope}[xshift=21cm]
\fill[blue!12] (-2.5,0) rectangle (2.5,1); \draw[thick] (-2.5,0)--(2.5,0) (-2.5,1)--(2.5,1);
\draw[ultra thick,red!70!black,dashed] (0,0)--(0,1);
\draw[->,very thick,green!50!black] (0,0.5)--(1.5,0.5); \draw[->,very thick,green!50!black] (0,0.5)--(-1.5,0.5);
\draw[->] (-2.8,0)--(2.8,0); \draw[->] (0,-0.3)--(0,4.8);
\node[font=\small] at (1.2,2) {$P_4$};
\end{scope}
\end{tikzpicture}
\end{center}
In rosso la parte $\mathrm{conv}\{v^i\}$ (segmento o vertici), in verde i generatori del cono $\mathrm{cono}\{e^j\}$. $P_1$: ogni punto è un punto del segmento tra $(1,2)$ e $(2,1)$ più una combinazione non negativa di $(1,0)$ e $(0,1)$. $P_2$: solo la parte convessa. $P_3$: segmento più una direzione. $P_4$: il segmento tratteggiato non è fatto di vertici, perché la striscia non ne ha.

\textbf{Come trovare la decomposizione in $\R^2$}: (1) disegna il poliedro; (2) segna i vertici: sono i $v^i$; (3) per ogni lato illimitato prendi la sua direzione: sono gli $e^j$ (equivalentemente risolvi $Ad\le0$ e prendi i raggi estremi del cono che ottieni).

Un poliedro ha almeno un vertice se e solo se non contiene rette. Ogni poliedro della forma $\{x: Ax\le b, x\ge0\}$ non vuoto (quindi ogni PL in forma canonica o standard) ha vertici.

\section{Proprietà della soluzione ottima}

\begin{prop}[caratterizzazione elementare dell'ottimalità]
Dato $(P)\ \max\{c^Tx: x\in P\}$ con $P=\{x\in\R^n:Ax\le b\}$:
\begin{enumerate}[nosep]
\item se $c\ne0$, nessuna soluzione ottima è un punto interno di $P$;
\item se $(P)$ ha due soluzioni ottime, allora ne ha infinite;
\item ogni soluzione ottima locale è anche ottima globale.
\end{enumerate}
\end{prop}
\begin{dimo}
\emph{1.} Se $x^*$ è interno, esiste $\varepsilon>0$ tale che $x^*+\varepsilon c\in P$. Ma $c^T(x^*+\varepsilon c)=c^Tx^*+\varepsilon\|c\|^2>c^Tx^*$: $x^*$ non è ottima.
\begin{center}
\begin{tikzpicture}[scale=0.55,>=Stealth]
\fill[blue!12] (0,0)--(5,0)--(5,2)--(3,4)--(0,3)--cycle; \draw[thick] (0,0)--(5,0)--(5,2)--(3,4)--(0,3)--cycle;
\draw[dashed,gray] (0.4,2.8)--(3.6,0.2); \draw[dashed,gray] (1.2,3.6)--(4.8,0.6);
\fill (2,1.5) circle (3pt) node[below left,font=\scriptsize] {$x^*$};
\draw[->,thick,red!70!black] (2,1.5)--(2.9,2.6) node[above,font=\scriptsize] {$x^*+\varepsilon c$};
\draw[gray] (2,1.5) circle (0.7);
\end{tikzpicture}
\end{center}
Attorno a un punto interno c'è una pallina tutta contenuta in $P$: muovendosi di poco lungo $c$ si resta ammissibili e si migliora.
\emph{2.} Se $x',x''$ sono ottime con valore $z^*$, per ogni $\lambda\in[0,1]$ il punto $\lambda x'+(1-\lambda)x''$ sta in $P$ (convesso) e ha valore $\lambda z^*+(1-\lambda)z^*=z^*$ (linearità): tutto il segmento è ottimo.
\emph{3.} Sia $\bar x$ ottimo locale e supponiamo esista $y\in P$ con $c^Ty>c^T\bar x$. I punti $\bar x+\lambda(y-\bar x)$, $\lambda\in(0,1]$, stanno in $P$ e hanno valore $c^T\bar x+\lambda(c^Ty-c^T\bar x)>c^T\bar x$; per $\lambda$ piccolo sono vicini quanto si vuole a $\bar x$, contro la sua ottimalità locale. \hfill$\square$
\end{dimo}
Le proprietà 2 e 3 valgono per qualunque problema con regione convessa e obiettivo lineare (la 3 anche con obiettivo concavo da massimizzare); sono la ragione per cui la PL è ``facile'' rispetto alla PLI.

\begin{prop}[Teorema fondamentale della PL]
Sia $P=\mathrm{conv}\{v^1,\dots,v^m\}+\mathrm{cono}\{e^1,\dots,e^p\}$. Se il problema $(P)\ \max\{c^Tx: x\in P\}$ ha valore ottimo finito, allora esiste $k\in\{1,\dots,m\}$ tale che $v^k$ è una soluzione ottima.
\end{prop}
\begin{dimo}
\emph{Passo 1.} $c^Te^j\le0$ per ogni $j$. Altrimenti, preso un $\bar x\in P$, i punti $\bar x+te^j$ ($t\ge0$) stanno in $P$ e $c^T(\bar x+te^j)=c^T\bar x+t\,c^Te^j\to+\infty$: il problema sarebbe illimitato.
\emph{Passo 2.} Sia $k$ tale che $c^Tv^k=\max_ic^Tv^i$. Ogni $x\in P$ è $x=\sum_i\lambda_iv^i+\sum_j\mu_je^j$ con $\lambda\ge0$, $\sum\lambda_i=1$, $\mu\ge0$, quindi
\[ c^Tx=\sum_i\lambda_i\,c^Tv^i+\sum_j\mu_j\,c^Te^j\ \le\ \sum_i\lambda_i\,c^Tv^i\ \le\ \Big(\sum_i\lambda_i\Big)c^Tv^k=c^Tv^k. \]
Poiché $v^k\in P$, $v^k$ è ottimo. \hfill$\square$
\end{dimo}

\begin{oss}[conseguenze]
\begin{itemize}[nosep]
\item Se una PL (con vertici) ha ottimo finito, \textbf{esiste un vertice ottimo}. Basta cercare tra i vertici, che sono in numero finito: al più $\binom{m}{n}$ (scelte di $n$ vincoli attivi tra $m$).
\item Enumerare tutti i vertici è però esponenziale. Il \textbf{simplesso} si muove invece da un vertice a uno \emph{adiacente} che migliora l'obiettivo, e si ferma quando nessun vertice adiacente migliora (per la proprietà 3, un ottimo ``locale sui vertici'' è globale).
\item Il teorema dice che \emph{esiste} un vertice ottimo, non che tutti gli ottimi sono vertici (caso (b): anche i punti interni allo spigolo sono ottimi).
\item Con il Passo 1: una PL ammissibile è illimitata se e solo se esiste una direzione di recessione $e$ con $c^Te>0$.
\end{itemize}
\end{oss}

\begin{esame}
Da saper fare: portare un problema in forma standard e canonica; risolvere graficamente un problema in 2 variabili (anche di minimo) e riconoscere i casi ottimo unico / infiniti ottimi / illimitato / inammissibile; verificare se un punto è combinazione convessa o conica di altri; dire se un punto è un vertice; trovare le direzioni di recessione ($Ad\le0$); scrivere la decomposizione $\mathrm{conv}+\mathrm{cono}$ di un poliedro in $\R^2$; enunciare e dimostrare la caratterizzazione elementare e il teorema fondamentale.
\end{esame}
